\documentclass[reqno]{exam}

\usepackage{amssymb, amsmath, amsthm, stmaryrd}
\usepackage{fullpage, parskip}
\usepackage{graphicx}
\usepackage{enumitem}

% Improved spacing
\usepackage[top=1in, bottom=0.8in, left=1in, right=1in]{geometry}


\title{\normalsize{
  JHED: \underline{\hspace{4cm}} \hspace{1em} Name: \underline{\hspace{9cm}}}
   }
\author{}
\date{}

% Custom spacing for better page utilization
\renewcommand{\baselinestretch}{0.95}

\begin{document}
\rhead{Computational Math, Spring 2026 - Practice Quiz}
\maketitle
\vspace{-2cm} % Reduce space after title

\begin{questions}

  \question
  Let $r_1$ and $r_2$ be the roots of the polynomial $p(x) = x^2 + (2 - \epsilon)x + 1$. Compute the condition number $\kappa_{r_1}(\epsilon)$. Simplify the denominator so that it's expressed entirely in terms of $r_1$ and $r_2$.

  \begin{solution}
    The polynomial is $p(x) = x^2 + (2 - \epsilon)x + 1$.

    Differentiating $r_1^2 + (2 - \epsilon)r_1 + 1 = 0$ with respect to $\epsilon$:
    $$2r_1\frac{\partial r_1}{\partial \epsilon} - r_1 + (2 - \epsilon)\frac{\partial r_1}{\partial \epsilon} = 0$$
    $$\frac{\partial r_1}{\partial \epsilon} = \frac{r_1}{2r_1 + (2 - \epsilon)}$$

    Using quadratic formula,
    \begin{align*}
      r_1                                             & = \frac{-(2 - \epsilon) + \sqrt{(2 - \epsilon)^2 - 4}}{2} \\
      r_2                                             & = \frac{-(2 - \epsilon) - \sqrt{(2 - \epsilon)^2 - 4}}{2} \\
      \implies
      (2 - \epsilon)                                  & = -(r_1 + r_2)                                            \\
      \implies \frac{\partial r_1}{\partial \epsilon} & = \frac{r_1}{r_1 - r_2}.
    \end{align*}


    Therefore:
    $$\kappa_{r_1}(\epsilon) = \left|\frac{\epsilon}{r_1} \cdot \frac{\partial r_1}{\partial \epsilon}\right| = \left|\frac{\epsilon}{r_1} \cdot \frac{r_1}{r_1-r_2}\right| = \left|\frac{\epsilon}{r_1-r_2}\right|.$$
  \end{solution}

  \newpage
  \question Let $f(x) = 1 - \cos x$. Find the condition number $\kappa_f(x)$. Determine the behavior of $\kappa_f(x)$ as $x \to 0$.
  \begin{solution}
    For $f(x) = 1 - \cos x$ with $f'(x) = \sin x$:
    $$\kappa_f(x) = \left|\frac{x \sin x}{1 - \cos x}\right|$$

    \textbf{As $x \to 0$:} We have the indeterminate form $\frac{0}{0}$. Applying L'Hôpital's rule:
    $$\lim_{x \to 0} \frac{x \sin x}{1 - \cos x} = \lim_{x \to 0} \frac{\sin x + x \cos x}{\sin x}$$

    This is still $\frac{0}{0}$. Applying L'Hôpital's rule again:
    $$\lim_{x \to 0} \frac{\sin x + x \cos x}{\sin x} = \lim_{x \to 0} \frac{\cos x + \cos x - x \sin x}{\cos x} = \frac{2}{1} = 2$$

    \textbf{Conclusion:} $\kappa_f(x) \to 2$ as $x \to 0$ (finite).
    
    \textbf{Note:} It is also possible to directly compute the limit without L'Hôpital's rule by using the Taylor expansion of $\cos x$ around $0$: $\cos x = 1 - \frac{x^2}{2} + \cdots$. Substituting this into the expression for $\kappa_f(x)$ and simplifying will yield the same result.
  \end{solution}

\end{questions}
\end{document}